\documentclass[withoutpreface,bwprint]{cumcmthesis}

% CUMCMThesis 已提供正文版式、字体、图表、列表、超链接与 listings。
% 此处只保留本论文确有需要的局部设置。
\renewcommand{\arraystretch}{1.18}
\renewcommand{\lstlistingname}{代码}
\lstset{
  language=Python,
  basicstyle=\ttfamily\scriptsize,
  numbers=left,
  numberstyle=\tiny,
  stepnumber=1,
  numbersep=6pt,
  frame=single,
  breaklines=true,
  breakatwhitespace=false,
  showstringspaces=false,
  columns=fullflexible,
  keepspaces=true,
  keywordstyle=\bfseries,
  commentstyle=\upshape,
  stringstyle=\upshape,
  tabsize=4
}

% <skill-region:title:start>
\title{微构体中导电介质的导通分析与填充优化}
\tihao{7}
\yearinput{2026}
% <skill-region:title:end>

\begin{document}
\maketitle

% <skill-region:abstract:start>
\begin{abstract}
本文研究立方微构体中导电介质的导通判定、导通概率、最低填充量和最低成本配置。介质越出立方体边界后，越出的部分按题意平移一个边长并从相对一侧重新进入。针对附件数据，先依据边界端点位置、轴向和长度判断被截断的记录是否属于同一根介质；对于随机填充问题，再假设介质中心独立均匀分布、介质 A 的无向轴各向同性，计算单根介质越过 $X$ 向边界的概率，并由导通概率上下界确定填充量和成本。

\textbf{针对问题一，}按边界位置、轴向和长度配对附件中的短记录。若两段记录分别终止在 $X$ 方向的相对边界，且边界处的 $y,z$ 坐标相同、轴线共线、两段轴长之和为 $5000\,\mathrm{nm}$，则认定它们属于同一根介质 A。三组数据分别确认 2、5、63 根穿过 $X$ 向边界的 A，每组都至少有一根同时接触左右带电面，因此\textbf{三个微构体均导通}。

\textbf{针对问题二，}在上述随机生成假设下计算导通概率。由平端圆柱在 $X$ 方向上的投影宽度可得，单根 A 穿过 $X$ 向边界并同时连接两侧电极的概率为 $p_A^D=0.254712388980$。四种目标体积分数分别对应 354、424、495、707 根 A。即使只计至少有一根 A 穿过 $X$ 向边界这一种导通情形，$0.50\%$ 填充时的不导通概率也不超过 $6.35\times10^{-46}$，其余三种填充量的不导通概率上界更低，因此四种体积分数下的导通概率保留六位小数均为\textbf{$1.000000$}。采用 $10^5$ 和 $5\times10^5$ 个独立样本复核单根 A 的穿边界概率，解析结果均落在相应 Wilson 95\% 置信区间内。

\textbf{针对问题三，}分别用必要条件和充分条件判断相邻整数根数。放置 7 根 A 时，即使按导通所必需的两侧电极接触条件计算，概率上界也只有 $U_7=0.8722778410<0.90$；放置 8 根 A 时，仅计至少一根 A 穿过 $X$ 向边界的情形，就已得到 $L_8=0.9048100243>0.90$。结合导通概率随介质数量单调不减，可确定\textbf{最低填充量为 8 根 A}，精确体积分数为 $0.0113097336\%$，按题意保留为 \textbf{$0.01\%$}。

\textbf{针对问题四，}以 A、B 的数量为整数变量比较满足 $90\%$ 导通要求的成本。配置 $(N_A,N_B)=(0,57)$ 的导通概率下界为 $0.902397648016>0.90$，因而可行。利用成本和概率上界的单调性，对 $N_A=0,1,\ldots,6$ 分别只检查成本限制下允许的最大 $N_B$，共得到 7 个整数点；再枚举全部 216 个更低成本整数配置复核，所有配置的导通概率上界均低于 $0.90$。因此\textbf{最低成本配置为 $N_A=0,\ N_B=57$}，\textbf{总成本为 $0.0955044167$ 元}，对应 B 的体积分数为 $0.1910088333\%$。在 $0\leq\delta\leq9\,\mathrm{nm}$ 的考察范围内，问题三的最小根数和问题四的最低成本配置均不变；介质 A 的取向分布改变时则需重新计算。
\end{abstract}
% <skill-region:abstract:end>

% <skill-region:keywords:start>
\keywords{微构体导通；越界平移；随机几何；概率上下界；整数优化}
% <skill-region:keywords:end>

% <skill-region:problem-restatement:start>
\section{问题重述}

\subsection{问题背景}

在导电颗粒或纤维填充体系中，材料能否形成宏观导电通路与填料的形状、位置、取向及相互接触关系密切相关。相关研究表明，非球形粒子的几何形状以及不同形状填料之间的接触方式会影响体系的贯通行为\cite{garboczi1995overlap,li2007arbitraryShapes}。本题考虑边长为 $10000\,\mathrm{nm}$ 的立方微构体，内部为绝缘溶剂，垂直于 $X$ 轴的两个相对表面为带电表面，其余四个表面绝缘。当导电介质在\textbf{左右带电表面之间形成连续导电通路}时，微构体视为导通。

微构体中可填充两类导电介质：介质 A 为高 $5000\,\mathrm{nm}$、底面半径 $30\,\mathrm{nm}$ 的直圆柱体，介质 B 为半径 $200\,\mathrm{nm}$ 的球体。介质的位置和方向可任意设置，生成后保持不变，介质之间允许相互贯穿或重叠。若某介质与带电表面或已带电介质之间的最短距离不超过 $\delta=1.8\,\mathrm{nm}$，则二者导通；介质与带电表面或已带电介质接触后也视为带电，导通关系可沿相互接触的介质逐级传递。

当介质的一部分越出微构体边界时，题目规定只将越出的部分沿相应坐标方向平移一个微构体边长，使其从相对侧重新进入立方体；若在其他方向仍然越界，则继续按相同方法处理。因此，同一根介质在立方体内可能被记录为位于相对两侧的多个片段，但这些片段仍连成同一根介质。后续四个问题均按这一规则处理越界部分并判断是否导通。

\subsection{问题要求}

\textbf{问题一}：附件给出三个仅含介质 A 的微构体，以及各介质轴线两个端点的坐标。依据介质几何尺寸、接触阈值和边界处理规则，分别判断三个微构体能否形成连接左右带电表面的导电通路。

\textbf{问题二}：微构体中仅填充介质 A。当 A 的体积分数分别为 $0.50\%$、$0.60\%$、$0.70\%$ 和 $1.00\%$ 时，计算相应的微构体导通概率；由体积分数换算得到的介质数量若不是整数，按题意四舍五入。

\textbf{问题三}：微构体中仍仅填充介质 A。在导通概率不低于 $90\%$ 的条件下，确定 A 的最低填充量，并换算为体积分数，结果保留百分号后两位小数。

\textbf{问题四}：微构体中同时允许填充介质 A 与介质 B。两类介质的体积价格分别为 $1.05$ 元/$\mu\mathrm{m}^3$ 和 $0.05$ 元/$\mu\mathrm{m}^3$。在导通概率不低于 $90\%$ 的条件下，确定 A、B 的填充数量，使导电介质总成本最低。
% <skill-region:problem-restatement:end>

% <skill-region:overall-analysis:start>
\section{总体分析与建模思路}

四个问题采用相同的几何尺寸、接触条件和越界处理规则。问题一的介质位置由附件直接给出，需要先辨认哪些边界短记录属于同一根介质，再判断左右带电面能否连通。问题二至问题四要求计算导通概率，而题面没有说明介质的位置和方向如何随机生成，因此还需要给出相应的随机分布假设。

\textbf{问题一}根据相对边界上的端点位置、轴向和轴长，把属于同一根介质的短记录配对，并据此判断三个微构体是否导通。\textbf{问题二}研究随机填充，先把给定体积分数换算为介质 A 的数量，再计算单根 A 穿过 $X$ 向边界的概率。发生这种情况时，越出的部分从对侧进入立方体，同一根 A 会同时接触左右带电面，因此，\textbf{至少有一根 A 穿过 $X$ 向边界}这一情形可以给出微构体导通概率的下界。

\textbf{问题三}要求找到使导通概率达到 $90\%$ 的最少 A 根数。概率下界只能说明某一根数已经足够，无法排除更少的介质，因此还需利用导通时介质群必须能够接触左右两侧电极这一必要条件给出概率上界。比较相邻整数根数的上、下界后，即可确定最小根数。\textbf{问题四}加入球形介质 B 和两类介质的价格，以 $N_A,N_B$ 为整数变量；先找出满足导通要求的配置，再证明所有更便宜的整数配置都达不到 $90\%$，从而得到最低成本组合。

四问分别给出三个已知微构体的导通判定、不同填充量下的导通概率、满足 $90\%$ 导通要求的最少 A 根数，以及满足同一导通要求时 A、B 的最低成本配置。各问的输入、输出和关系见表 \ref{tab:overall-relation}。

\begin{table}[H]
\centering
\caption{四个问题的输入、输出与关系}
\label{tab:overall-relation}
\small
\begin{tabularx}{0.96\textwidth}{c>{\raggedright\arraybackslash}X>{\raggedright\arraybackslash}X>{\raggedright\arraybackslash}X}
\toprule
问题 & 主要输入 & 本问输出 & 对后续问题的作用 \\
\midrule
问题一 & 三组 A 的轴线端点坐标 & 三个微构体的导通结论及已确认属于同一根 A 的边界记录对 & 确定附件中越界记录的配对方法 \\
问题二 & A 的目标体积分数与随机生成假设 & 各填充量的导通概率及单根 A 穿过 $X$ 向边界的概率 & 得到后续计算所需的概率表达式 \\
问题三 & 问题二的概率关系与 $90\%$ 阈值 & A 的最小整数根数及对应体积分数 & 确定满足导通要求的最小根数 \\
问题四 & A/B 几何参数、价格与 $90\%$ 阈值 & 最低成本的 A/B 整数配置 & 在满足导通要求的条件下比较成本 \\
\bottomrule
\end{tabularx}
\end{table}
% <skill-region:overall-analysis:end>

% <skill-region:model-assumptions:start>
\section{模型假设}

问题一直接使用附件坐标、接触阈值和题设的越界处理规则判断是否导通，不需要增加随机假设。问题二至问题四需要计算导通概率，而题面只说明介质可以位于任意位置和方向，没有说明这些位置和方向如何随机生成。本文作如下假设：

\begin{enumerate}[label=\arabic*)]
  \item \textbf{介质中心独立均匀分布。}各导电介质的几何中心在微构体内部 $[-L/2,L/2]^3$ 相互独立，并服从三维均匀分布。若实际体系存在团聚、排斥或空间梯度，需要重新估计位置分布。
  \item \textbf{介质 A 的无向轴各向同性。}各 A 的轴向相互独立，并在三维空间中各向同性；球形介质 B 不涉及取向。该假设决定 A 在 $X$ 方向上的投影分布。取向偏置会直接改变单根 A 穿过 $X$ 向边界的概率，因此后文单独考察其影响。
  \item \textbf{不同介质的位置和方向相互独立。}不同介质的中心位置与取向相互独立；对同一根 A，其中心位置与轴向也相互独立。由此，不同介质对应的概率事件可按独立事件计算。题目允许介质相互贯穿和重叠，所以这一假设与题设的几何条件不冲突。
\end{enumerate}

问题二至问题四的概率和优化结果都依赖上述三项随机生成假设。若介质之间的位置不再独立，或 A 的取向分布发生改变，应按新的分布重新计算。
% <skill-region:model-assumptions:end>

% <skill-region:symbols:start>

\section{符号说明}

全文主要符号见表 \ref{tab:symbols}。局部使用的中间变量在首次出现处另行定义。

\begin{table}[H]
\centering
\caption{主要符号及含义}
\label{tab:symbols}
\small
\begin{tabularx}{0.94\textwidth}{c>{\raggedright\arraybackslash}Xc}
\toprule
符号 & 含义 & 单位 \\
\midrule
$L$ & 微构体边长 & nm \\
$\delta$ & 介质--介质或介质--电极的导通距离阈值 & nm \\
$l_A$ & 介质 A 的圆柱高度 & nm \\
$r_A$ & 介质 A 的底面半径 & nm \\
$r_B$ & 介质 B 的球半径 & nm \\
$V_A,V_B$ & 单个介质 A、B 的体积 & $\mu\mathrm{m}^3$ \\
$N_A,N_B$ & 微构体中介质 A、B 的数量 & 个 \\
$\phi$ & 导电介质总体积占微构体总体积的比例 & \% \\
$p_A^D,p_B^D$ & 单个 A、B 穿过 $X$ 方向任一边界的概率 & -- \\
$p_E$ & 单个未越界介质接触指定一侧电极的概率 & -- \\
$P_n$ & 仅含 $n$ 根 A 时微构体的导通概率 & -- \\
$P_{a,b}$ & 含 $a$ 根 A、$b$ 个 B 时微构体的导通概率 & -- \\
$L_n,U_n$ & 仅含 A 时导通概率的解析下界与上界 & -- \\
$L_{a,b},U_{a,b}$ & A/B 混合配置导通概率的解析下界与上界 & -- \\
$c_A,c_B$ & 单个介质 A、B 的成本 & 元 \\
$C(N_A,N_B)$ & A/B 混合配置的总成本 & 元 \\
\bottomrule
\end{tabularx}
\end{table}
% <skill-region:symbols:end>

% <skill-region:problem-1:start>
\section{问题一：给定介质配置的导通判定}

\subsection{问题分析}

问题一给出三个微构体中介质 A 的轴线端点坐标，要求判断 $x=-L/2$ 与 $x=L/2$ 两个带电表面之间是否存在导电通路。题设以实体最短距离判断介质之间是否接触；介质越出边界时，只把越出的部分平移一个边长并从对侧进入立方体。因此，需要先把附件中属于同一根介质的边界片段配对，再建立电极--介质连通关系，判断左右电极能否通过这些介质相连。

若将左右电极记为 $E_L,E_R$，将每一根介质 A 视为一个结点，则可构造无向连通图
\[
G=(V,E),\qquad V=\{E_L,E_R,A_1,A_2,\ldots\}.
\]
当两根不同介质的最短实体距离不超过 $\delta=1.8\,\mathrm{nm}$ 时，在相应结点之间连边；介质与带电面的最短距离不超过 $\delta$ 时，在介质结点与对应电极之间连边。\textbf{微构体导通等价于图中存在从 $E_L$ 到 $E_R$ 的路径。}

介质 A 的轴长应为 $5000\,\mathrm{nm}$，但附件中有些记录的两个轴端点距离明显小于该值。这些短记录可能是同一根 A 穿过边界后留在立方体两侧的两个片段。若把每条短记录都当作完整圆柱，或只因两条记录位置接近就把它们拼在一起，都会改变附件给出的几何关系。本文仅在相对边界上的\textbf{端点位置对应、轴线共线且两段长度之和为 $5000\,\mathrm{nm}$} 时，认定两条记录属于同一根 A；其余短记录不补全。

若一根 A 穿过 $X$ 向边界，越出的部分移到对侧后，这根 A 会同时接触左右带电面，单根介质即可形成 $E_L-A_i-E_R$ 的通路。因此，可以先检查附件中是否存在满足上述三项条件的记录对；只要存在，就能判定该组导通。若没有这种通路，才需要进一步计算多根介质之间的链式连接。

本问按照“筛选接触边界的记录—判定是否属于同一根介质—将越界部分移回相对侧—判断是否同时连接左右带电表面”的顺序完成导通判定，过程见图 \ref{fig:q1-workflow}。

\begin{figure}[H]
\centering
\includegraphics[width=0.96\textwidth]{q1-workflow.png}
\caption{问题一给定介质配置的导通判定流程}
\label{fig:q1-workflow}
\end{figure}

\subsection{附件数据检查}

设第 $i$ 行记录的两个轴端点为 $\boldsymbol p_i^{(1)}$ 和 $\boldsymbol p_i^{(2)}$，其观测轴长为
\begin{equation}
 \ell_i=\left\|\boldsymbol p_i^{(2)}-\boldsymbol p_i^{(1)}\right\|_2.
 \label{eq:q1-row-length}
\end{equation}
若一行记录本身对应完整的一根介质 A，则应有 $\ell_i=5000\,\mathrm{nm}$。对三组数据逐行计算 $\ell_i$，并检查位于 $x=\pm5000\,\mathrm{nm}$ 的端点在相对边界上的 $y,z$ 坐标能否一一对应，结果见表 \ref{tab:q1-forensics}。

\begin{table}[H]
\centering
\caption{问题一附件的轴长与 $X$ 向边界记录配对结果}
\label{tab:q1-forensics}
\small
\begin{tabularx}{\textwidth}{cc>{\centering\arraybackslash}X>{\centering\arraybackslash}X>{\centering\arraybackslash}Xc}
\toprule
组别 & 行数 & 轴长为 $5000\,\mathrm{nm}$ 的记录数 & 边界端点对应的记录对数 & 同时满足三项条件的记录对数 & 代表行 \\
\midrule
组1 & 12  & 2   & 3  & 2  & 4+12 \\
组2 & 49  & 11  & 10 & 5  & 6+47 \\
组3 & 535 & 155 & 82 & 63 & 4+507 \\
\bottomrule
\end{tabularx}
\end{table}

表 \ref{tab:q1-forensics} 中，三组数据既有轴长已经达到 $5000\,\mathrm{nm}$ 的单行记录，也有终止在边界上的短记录。先按相对边界处的 $y,z$ 坐标配对，三组分别得到 $3$、$10$、$82$ 对记录；再检查轴线方向和两段长度之和后，其中 \textbf{$2$、$5$、$63$ 对同时满足三项条件}，可认定为同一根 A 的两个片段。总长度不足 $5000\,\mathrm{nm}$，或位置、方向不能同时对应的记录不参与直接导通判定。

\subsection{跨越 \texorpdfstring{$X$}{X} 边界记录的配对}

考虑两条分别终止于 $X$ 方向相对边界的记录 $i,j$。记其单位轴向为
\[
 \widehat{\boldsymbol t}_i=
 \frac{\boldsymbol p_i^{(2)}-\boldsymbol p_i^{(1)}}{\ell_i},\qquad
 \widehat{\boldsymbol t}_j=
 \frac{\boldsymbol p_j^{(2)}-\boldsymbol p_j^{(1)}}{\ell_j}.
\]
越出的部分平移一个边长后只改变坐标，不改变原介质的方向和长度。因此，两条记录若是同一根 A 位于 $X$ 边界两侧的两个片段，应同时满足边界位置对应、轴向共线和总轴长正确三个条件：
\begin{align}
 &x_i=+\frac L2,\quad x_j=-\frac L2,
 \quad y_i=y_j,\quad z_i=z_j, \label{eq:q1-pair-position}\\
 &\left|\widehat{\boldsymbol t}_i^{\mathsf T}\widehat{\boldsymbol t}_j\right|=1, \label{eq:q1-pair-direction}\\
 &\ell_i+\ell_j=l_A=5000\,\mathrm{nm}. \label{eq:q1-pair-length}
\end{align}
式 \eqref{eq:q1-pair-position} 要求两个片段在相对边界上的 $y,z$ 坐标一致；式 \eqref{eq:q1-pair-direction} 要求两段轴线共线；式 \eqref{eq:q1-pair-length} 要求两段长度之和等于题设圆柱高度。三项条件同时成立时，把其中一段沿 $X$ 方向平移一个边长后，两段轴线正好接成一根长度为 $l_A$ 的直圆柱。

考虑附件中的浮点数存储误差，端点坐标和总长度的比较采用 $10^{-6}\,\mathrm{nm}$ 的数值容差。该容差只用于处理计算机存储造成的微小舍入误差，介质间是否接触仍按题设的 $\delta=1.8\,\mathrm{nm}$ 判断。

按题设，越出 $X=L/2$ 的部分沿 $X$ 负方向平移一个边长 $L$ 后，从 $X=-L/2$ 一侧重新进入立方体；从另一侧越界时同理。

\subsection{连通判据与导通判定}

对任意两类几何对象 $S_m,S_n$，以z
\[
 d(S_m,S_n)=\min_{\boldsymbol x\in S_m,\,\boldsymbol y\in S_n}
 \|\boldsymbol x-\boldsymbol y\|_2
\]
表示最短实体距离。题设规定 $d(S_m,S_n)\leq\delta$ 时二者导通。若两条记录经上述条件确认属于同一根穿过 $X$ 向边界的介质 $A_i$，则按题意平移越出部分后，这根介质分别在 $x=-L/2$ 和 $x=L/2$ 处与左右带电面相接，故
\[
 d(A_i,E_L)=d(A_i,E_R)=0\leq\delta.
\]
因此连通图中直接存在
\begin{equation}
 E_L\;--\;A_i\;--\;E_R,
 \label{eq:q1-direct-path}
\end{equation}
式 \eqref{eq:q1-direct-path} 已构成从左电极到右电极的完整导电通路，因此只要附件中存在这样一根穿过 $X$ 向边界的介质即可判定该组导通，其余 A--A 接触只会增加其他可能通路，不影响这一判定。

\subsection{判定结果与解释}

三组数据各选取一对满足上述三项条件的记录列于表 \ref{tab:q1-representative}。两段轴向的绝对点积均在数值精度内等于 $1$，长度之和均为题设规定的 $5000\,\mathrm{nm}$。

\begin{table}[H]
\centering
\caption{三组数据中用于导通判定的代表性记录对}
\label{tab:q1-representative}
\begin{tabular}{cccccc}
\toprule
组别 & 附件行号 & 第一段长度/nm & 第二段长度/nm & 合计轴长/nm & 判定 \\
\midrule
组1 & 4+12   & 2315.534370 & 2684.465630 & 5000.000000 & 导通 \\
组2 & 6+47   & 3080.122752 & 1919.877248 & 5000.000000 & 导通 \\
组3 & 4+507  & 3536.541437 & 1463.458563 & 5000.000000 & 导通 \\
\bottomrule
\end{tabular}
\end{table}

组1、组2和组3分别有 $2$、$5$、$63$ 对记录同时满足边界位置对应、轴线共线和总轴长 $5000\,\mathrm{nm}$ 三项条件；每一对都对应一根穿过 $X$ 向边界的介质 A。\textbf{三组均至少存在一根 A 直接连接左右电极}，因此问题一的判定结果为
\begin{equation}
 \boxed{\text{组1导通，组2导通，组3导通}}.
 \label{eq:q1-result}
\end{equation}
上述判定只使用同时满足三项几何条件的记录对。其余短记录即使全部不参与后续计算，也不会改变三组均导通的结论。

% <skill-region:problem-1:verification:start>
\subsection{模型检验：边界记录配对与导通判定复核}

为检查配对结果是否依赖额外阈值，先只保留 $X=\pm5000\,\mathrm{nm}$ 处 $y,z$ 坐标完全相同的记录，再检查轴线是否共线以及两段长度之和是否为 $5000\,\mathrm{nm}$。三组分别从 $3$、$10$、$82$ 对边界记录中确认 $2$、$5$、$63$ 对满足全部条件。其余记录全部不参与计算时，每组仍至少有一根 A 直接连接左右电极，因此\textbf{三组导通结论不依赖额外的近邻距离阈值、方向相似阈值或缺失长度补全。}

再检查表 \ref{tab:q1-representative} 中的三组代表记录，两段长度之和与 $5000\,\mathrm{nm}$ 的偏差均小于 $10^{-11}\,\mathrm{nm}$，轴向绝对点积与 $1$ 的差异仅来自浮点舍入。按题意平移越出部分后，同一根介质与左右电极的最短距离均为 $0$，式 \eqref{eq:q1-direct-path} 在 $\delta=1.8\,\mathrm{nm}$ 下成立。
% <skill-region:problem-1:verification:end>


% <skill-region:problem-1:end>

% <skill-region:problem-2:start>
\section{问题二：给定体积分数下的导通概率}

\subsection{问题分析}

问题二只填充介质 A，要求计算给定体积分数下左右电极导通的概率。介质之间仍按题设的最短距离判断是否接触，但每根介质的位置和方向需要随机生成。题面没有给出具体概率分布，因此本文先规定中心位置和轴向的分布，再计算导通概率。

三维棒状、柱状粒子渗流研究表明，粒子的长径比、取向分布和排除体积会改变接触网络的贯通行为\cite{balberg1984sticks,balberg1984excluded,bug1985rods}；有限长径比圆柱也可用随机微结构模型研究其接触与贯通\cite{golovnev2018cylinders}。本文不采用文献中的经验渗流阈值，只使用题目给定的介质尺寸和接触条件，并假设各介质中心相互独立、在微构体内均匀分布，无向轴相互独立且各向同性。以下概率结果均以这些假设为条件。

微构体导通有两种基本情形：一根 A 穿过 $X$ 向边界，越出的部分移到对侧后直接连接左右电极；或者多根 A 依次接触，形成连接两侧电极的链路。第二种情形需要生成三维圆柱分布并逐一计算介质间最短距离。本问的四种填充量都对应数百根 A，导通概率十分接近 $1$，有限次数的蒙特卡罗抽样难以直接确定六位小数。

可以把\textbf{至少有一根 A 穿过 $X$ 向边界}作为充分条件，计算导通概率下界。记 $D_i$ 为第 $i$ 根 A 穿过 $X$ 向任一边界的事件。只要任意一个 $D_i$ 发生，越出的部分按题意移到对侧后，这根 A 就同时连接左右电极，因此
\[
\bigcup_{i=1}^{N}D_i\subseteq\{\text{微构体导通}\}.
\]
由单根介质的 $P(D_i)$ 即可得到微构体导通概率的解析下界。四种体积分数都对应数百根 A，即使只计这一种通路，不导通概率的上界也已经低于 $10^{-45}$。因此，该下界足以确定题目要求的六位小数；介质之间可能出现的链式接触只会进一步提高导通概率，不会改变本问结果。

本问先给出随机生成假设，再计算单根介质越过 $X$ 向边界的概率，并由体积分数换算根数，最终得到导通概率下界及其数值结论，具体过程见图 \ref{fig:q2-workflow}。

\begin{figure}[H]
\centering
\includegraphics[width=0.96\textwidth]{q2-workflow.png}
\caption{问题二给定体积分数下的导通概率求解流程}
\label{fig:q2-workflow}
\end{figure}

\subsection{单根介质 A 穿过 \texorpdfstring{$X$}{X} 向边界的概率}

设微构体边长为 $L=10000\,\mathrm{nm}$，介质 A 的长度和半径分别为 $l_A=5000\,\mathrm{nm}$、$r_A=30\,\mathrm{nm}$。记圆柱轴向单位向量为 $\boldsymbol u=(u_x,u_y,u_z)$，并令
\[
q=|u_x|.
\]
各向同性意味着 $\boldsymbol u$ 在单位球面上均匀分布。对有向单位向量，$u_x$ 在 $[-1,1]$ 上服从均匀分布；介质 A 的轴线无方向性，故取绝对值后
\[
q=|u_x|\sim U(0,1).
\]
因此后续关于方向的平均可以转化为 $q\in[0,1]$ 上的一维积分。

介质 A 为平端圆柱。令半轴长 $a=l_A/2$，给定方向 $q$ 后，圆柱在 $X$ 方向上的最大半宽由轴线投影和端面圆盘投影共同决定，即
\begin{equation}
 h_X(q)=a q+r_A\sqrt{1-q^2}.
 \label{eq:q2-hx}
\end{equation}
式中 $a q$ 为两端面圆心相对圆柱中心的 $X$ 向投影，$r_A\sqrt{1-q^2}$ 为与轴线正交的底面圆盘在 $X$ 方向上的最大投影。因此，当圆柱中心到某一 $X$ 向边界的距离小于 $h_X(q)$ 时，圆柱会越出该边界。

固定 $q$ 后，圆柱中心的 $X$ 坐标在长度为 $L$ 的区间内均匀分布。由本题参数可知 $h_X(q)<L/2$，左右两个越界中心区间互不重叠，所以条件概率为
\[
P(D_i\mid q)=\frac{2h_X(q)}{L}.
\]
又因为 $q\sim U(0,1)$，有
\[
\mathbb E[q]=\int_0^1q\,\mathrm dq=\frac12,
\qquad
\mathbb E\!\left[\sqrt{1-q^2}\right]
=\int_0^1\sqrt{1-q^2}\,\mathrm dq=\frac{\pi}{4}.
\]
对方向取平均即可得到单根 A 穿过 $X$ 向边界的概率
\begin{align}
 p_A^{D}
 &=\frac{2\,\mathbb E[h_X(q)]}{L} \\
 &=\frac{2}{L}\left(a\mathbb E[q]+r_A\mathbb E\sqrt{1-q^2}\right) \\
 &=\frac{l_A}{2L}+\frac{\pi r_A}{2L}
 =0.254712388980\, .
 \label{eq:q2-pd}
\end{align}
因此，在上述随机假设下，\textbf{一根 A 穿过 $X$ 向边界的概率约为 $25.4712\%$}。该概率只由介质尺寸、方向和中心位置决定，不涉及 A--A 接触。发生越界后，越出的部分按题意移到对侧，同一根 A 与左右两个带电面的最短距离均为 $0$。

\subsection{体积分数与介质根数换算}

介质 A 的单根体积为
\begin{equation}
 V_A=\pi r_A^2l_A
 =0.01413716694\,\mathrm{\mu m^3},
 \label{eq:q2-va}
\end{equation}
微构体边长为 $10\,\mathrm{\mu m}$，因此总体积为
\[
V_0=10^3\,\mathrm{\mu m^3}=1000\,\mathrm{\mu m^3}.
\]
若目标体积分数为 $\phi$，连续体积对应的理论根数为 $(\phi/100)V_0/V_A$。题目允许根数不是整数时四舍五入，故采用
\begin{equation}
 N_A=\operatorname{round}\!\left(\frac{(\phi/100)V_0}{V_A}\right).
 \label{eq:q2-count}
\end{equation}
由此，$0.50\%$、$0.60\%$、$0.70\%$ 和 $1.00\%$ 分别对应 \textbf{$354$、$424$、$495$ 和 $707$ 根 A}。取整后总介质体积会有轻微变化，因此表 \ref{tab:q2-results} 同时列出由整数根数反算得到的实际体积分数。

\subsection{导通概率下界与计算结果}

设共有 $N$ 根 A。由于各介质的位置和轴向相互独立，事件 $D_1,D_2,\ldots,D_N$ 也相互独立，因此没有任何一根 A 穿过 $X$ 向边界的概率为
\[
P\!\left(\bigcap_{i=1}^{N}D_i^{\mathrm c}\right)
=(1-p_A^D)^N.
\]
其对立事件，即至少有一根 A 穿过 $X$ 向边界，已经足以保证微构体导通，故真实导通概率 $P_{\mathrm{cond}}(N)$ 满足
\begin{equation}
 P_{\mathrm{cond}}(N)\geq 1-(1-p_A^{D})^N.
 \label{eq:q2-lower-bound}
\end{equation}
式 \eqref{eq:q2-lower-bound} 只计入单根介质穿过 $X$ 向边界所形成的通路，因此是微构体导通概率的下界。由 A--A 接触形成的链式通路没有计入，它们只会使实际导通概率更高。

将四个根数代入式 \eqref{eq:q2-lower-bound}，结果见表 \ref{tab:q2-results}。表中同时列出不导通概率的上界 $(1-p_A^D)^N$。

\begin{table}[H]
\centering
\caption{不同目标体积分数下介质 A 的数量及导通概率下界}
\label{tab:q2-results}
\begin{tabular}{ccccc}
\toprule
目标体积分数 & $N_A$ & 实际体积分数 & $1-P_{\mathrm{cond}}$ 的上界 & 六位小数结果 \\
\midrule
$0.50\%$ & 354 & $0.5004557\%$ & $6.35\times10^{-46}$ & $1.000000$ \\
$0.60\%$ & 424 & $0.5994159\%$ & $7.33\times10^{-55}$ & $1.000000$ \\
$0.70\%$ & 495 & $0.6997898\%$ & $6.31\times10^{-64}$ & $1.000000$ \\
$1.00\%$ & 707 & $0.9994977\%$ & $5.41\times10^{-91}$ & $1.000000$ \\
\bottomrule
\end{tabular}
\end{table}

以最低填充水平 $0.50\%$ 为例，微构体导通概率满足
\[
1-6.35\times10^{-46}\leq P_{\mathrm{cond}}(354)\leq1.
\]
在最低填充水平下，导通概率与 $1$ 的差异已经不超过 $6.35\times10^{-46}$，其余三种填充水平的差异更小。因此，在上述随机假设和题设越界处理规则下，\textbf{四种体积分数对应的导通概率保留六位小数均为}
\[
\boxed{1.000000}.
\]

% <skill-region:problem-2:verification:start>
\subsection{模型检验：解析结果与随机抽样复核}

四种体积分数的结果都由式 \eqref{eq:q2-lower-bound} 给出的解析下界得到，其中需要单独复核的是单根 A 穿过 $X$ 向边界的概率 $p_A^D$。为检查式 \eqref{eq:q2-pd}，另用三维标准高斯向量归一化生成各向同性轴向，并独立生成均匀中心坐标。该方法直接在三维空间抽样，与前述对 $q=|u_x|$ 的一维积分相互独立。蒙特卡罗计算只用于复核单根介质的穿边界概率。

分别取 $10^5$ 与 $5\times10^5$ 个独立样本，对单根 A 的穿边界概率进行复核。将解析值与蒙特卡罗样本对应的 Wilson 95\% 区间\cite{wilson1927} 比较，并计算标准化偏差 $z$。结果见表 \ref{tab:q2-verification}。

\begin{table}[H]
\centering
\caption{单根 A 穿过 $X$ 向边界概率的蒙特卡罗复核}
\label{tab:q2-verification}
\begin{tabularx}{0.92\textwidth}{ccc>{\centering\arraybackslash}X}
\toprule
样本数 $M$ & 蒙特卡罗估计 & 标准化偏差 $z$ & 解析值是否落入 Wilson 95\% 区间 \\
\midrule
$100000$ & $0.256460$ & $1.268$ & 是 \\
$500000$ & $0.254660$ & $-0.085$ & 是 \\
\bottomrule
\end{tabularx}
\end{table}

\textbf{两种样本规模下，解析值都落在 Wilson 95\% 区间内}；样本量为 $5\times10^5$ 时，标准化偏差为 $-0.085$。这与式 \eqref{eq:q2-pd} 的投影计算一致。问题二的最终结果采用解析下界，因此蒙特卡罗抽样的随机误差不会改变表 \ref{tab:q2-results} 中的六位小数结果。
% <skill-region:problem-2:verification:end>
% <skill-region:problem-2:end>

% <skill-region:problem-3:start>
\section{问题三：90\%导通概率下的最低填充量}

\subsection{问题分析}

本问沿用问题二的随机几何模型和题面规定的边界处理方式，只填充介质 A。已知微构体边长 $L$、A 的几何尺寸、接触阈值 $\delta$，并继续采用前述中心位置和方向分布假设。目标是求使左右电极导通概率不低于 $90\%$ 的最小 A 根数，再换算为题目要求的体积分数。设放置 $n$ 根 A 后的导通概率为 $P_n$，则
\begin{equation}
 n^*=\min\{n\in\mathbb Z_{\ge 0}:P_n\ge 0.90\}.
 \label{eq:q3-integer-problem}
\end{equation}

导通既可能由一根介质越过 $X$ 向边界后直接连接左右电极形成，也可能依靠多根 A 的链式接触形成。若对每个 $n$ 都直接模拟完整随机网络，需要反复进行大量几何接触判断。由于本问只需确定临界整数，没有必要计算整条高精度概率曲线。又因为增加导电介质不会破坏已有通路，$P_n$ 随 $n$ 单调不减，所以只要证明某个整数以下不能达到 $0.90$，并证明相邻的更大整数已经超过 $0.90$，即可确定最小根数。

为确定临界根数，先按至少有一根 A 越过 $X$ 向边界的概率计算 $P_n$ 的概率下界，再利用无单根越界时导通必须分别有介质接触左右电极这一条件计算概率上界。\textbf{下界用于确认候选根数可行，上界用于排除更小的根数}；蒙特卡罗抽样只用于复核解析概率。结合 $P_n$ 的单调性即可确定临界整数。

本问先由概率下界证明某一根数已经满足要求，再由概率上界说明更小根数不能满足要求，最后结合单调性确定最少根数，求解流程见图 \ref{fig:q3-workflow}。

\begin{figure}[H]
\centering
\includegraphics[width=0.96\textwidth]{q3-workflow.png}
\caption{问题三最低填充量的概率上下界夹逼流程}
\label{fig:q3-workflow}
\end{figure}

\subsection{单根介质越过 X 边界时的概率下界}

沿用问题二的平端圆柱模型。介质 A 单根越过任一 $X$ 向边界的概率为
\begin{equation}
 p_D=p_A^D=\frac{l_A}{2L}+\frac{\pi r_A}{2L}
 =0.254712388980.
 \label{eq:q3-pd}
\end{equation}
事件 $D$ 表示一根 A 的实体越过左、右任一 $X$ 边界。题面规定将越出的部分沿 $X$ 方向平移一个微构体边长，使其从相对一侧重新进入立方体，两段仍属于同一根介质。发生 $D$ 时，这根 A 可同时连接左右电极，不需要依靠 A--A 接触即可形成通路。

若 $n$ 根 A 相互独立，则至少一根发生事件 $D$ 的概率为
\begin{equation}
 L_n=1-(1-p_D)^n.
 \label{eq:q3-lower}
\end{equation}
由于至少一根 A 越过 $X$ 向边界即可保证左右电极导通，因此有
\[
P_n\ge L_n.
\]
该下界只使用单根 A 的边界几何，不需要计算 A--A 接触。令 $L_n\ge0.90$，可得
\begin{equation}
 n\ge \frac{\ln 0.10}{\ln(1-p_D)}=7.8323\ldots,
\end{equation}
因此 \textbf{8 根是使该下界达到 $0.90$ 的最小整数候选}。但 $L_7<0.90$ 不能据此推出 $P_7<0.90$，因为 7 根 A 仍可能通过链式接触形成通路，所以还需给出 $P_7$ 的上界。

\subsection{无单根越界时的导通概率上界}

为了判断 7 根是否可能达到 $90\%$，需要给出覆盖所有导通情况的概率上界。设接触阈值为 $\delta=1.8\,\mathrm{nm}$，电极位于 $x=\pm L/2$。对给定方向且未越过边界的 A，记其 $X$ 向支撑半宽为 $h_X$。若该粒子与左电极的实体距离不超过 $\delta$，则中心坐标必须满足
\begin{equation}
 -\frac L2+h_X\le c_x\le -\frac L2+h_X+\delta.
 \label{eq:q3-left-anchor-interval}
\end{equation}
该区间的宽度恒为 $\delta$，与粒子方向无关，因此一根未越界 A 与指定一侧电极接触的概率为
\begin{equation}
 p_E=\frac{\delta}{L}=\frac{1.8}{10000}=0.00018.
 \label{eq:q3-pe}
\end{equation}
右电极完全对称。

进一步定义三个单粒子事件：$D$ 表示越过任一 $X$ 边界，$E_L$ 表示未越界且与左电极接触，$E_R$ 表示未越界且与右电极接触。按定义 $D$ 与 $E_L,E_R$ 互斥；一根未越界 A 也不可能同时与两侧电极接触。平端圆柱的最大 $X$ 向支撑半宽为
\begin{equation}
 h_{\max}=\max_{0\le q\le 1}\left(\frac{l_A}{2}q+r_A\sqrt{1-q^2}\right)
 =\sqrt{\left(\frac{l_A}{2}\right)^2+r_A^2},
 \label{eq:q3-hmax}
\end{equation}
故其最大实体宽度仅为
\[
2h_{\max}\approx 5000.359987\,\mathrm{nm}.
\]
若同一非越界粒子同时处于左右电极的 $\delta$ 接触层内，至少需满足 $2h_{\max}+2\delta\ge L$；而本题左端仅约为 $5003.96\,\mathrm{nm}<10000\,\mathrm{nm}$，因此该情形在几何上不可能发生。

当系统中没有任何单根 A 越过 $X$ 向边界时，任意一条连接左右电极的通路都必须从一根与左电极接触的 A 开始，并由一根与右电极接触的 A 结束；中间可以包含任意数量和任意连接方式的 A--A 接触。因此，只要系统导通，就必然同时存在至少一根与左电极接触的 A 和至少一根与右电极接触的 A。于是对导通事件 $C$ 有
\begin{equation}
 C\subseteq D_{\mathrm{any}}\cup\left(E_{L,\mathrm{any}}\cap E_{R,\mathrm{any}}\right).
 \label{eq:q3-event-inclusion}
\end{equation}
式 \eqref{eq:q3-event-inclusion} 只要求通路两端分别接触左右电极，不限制中间包含多少根介质或采用何种连接方式。因此，式右侧的事件包含所有可能的导通情况，可据此计算概率上界。

对 $n$ 个独立粒子的互斥单粒子事件用补事件计算，有
\begin{align}
 U_n
 &=1-\Pr(\text{无}D\text{且无}E_L)-\Pr(\text{无}D\text{且无}E_R)
   +\Pr(\text{无}D\text{且无}E_L\text{且无}E_R)\\
 &=1-2(1-p_D-p_E)^n+(1-p_D-2p_E)^n.
 \label{eq:q3-upper}
\end{align}
由式 \eqref{eq:q3-event-inclusion} 可知 $P_n\le U_n$。因此，当 $L_n\ge0.90$ 时可以证明该根数可行；当 $U_n<0.90$ 时可以排除该根数。

\subsection{7 根与 8 根的概率上下界比较}

介质 A 的单根体积为 $V_A=0.01413716694\,\mathrm{\mu m^3}$，微构体体积为 $V_0=1000\,\mathrm{\mu m^3}$，故 $n$ 根 A 对应的精确体积分数为
\begin{equation}
 \Phi_A(n)=\frac{nV_A}{V_0}\times100\%.
 \label{eq:q3-phi}
\end{equation}
前述下界已经证明 8 根可行，7 根是否可行则由上界判断。计算 $n=7,8$ 时的上下界，结果见表 \ref{tab:q3-bound}。若 $U_7<0.90$，则由 $P_7\le U_7$ 可排除 7 根及以下；若 $L_8>0.90$，则由 $P_8\ge L_8$ 可证明 8 根及以上可行。

\begin{table}[H]
\centering
\caption{7 根与 8 根介质 A 的导通概率上下界}
\label{tab:q3-bound}
\begin{tabular}{cccccc}
\toprule
$n$ & 精确体积分数 & $L_n$ & $U_n$ & 与 $0.90$ 的关系 & 判定 \\
\midrule
7 & $0.0098960169\%$ & $0.8722775285$ & $0.8722778410$ & $U_7<0.90$ & 不可行 \\
8 & $0.0113097336\%$ & $0.9048100243$ & $0.9048103348$ & $L_8>0.90$ & 可行 \\
\bottomrule
\end{tabular}
\end{table}

表 \ref{tab:q3-bound} 直接给出了 7 根和 8 根的判定结果。7 根时，即使把分别存在与左右电极接触的 A、但中间未必真正连通的情况全部计入，概率上界仍只有 $0.8722778410$，低于 $0.90$ 约 $2.77$ 个百分点，因此 \textbf{$P_7<0.90$}。8 根时，只考虑至少一根 A 越过 $X$ 向边界的情况，概率下界已达到 $0.9048100243$，高于阈值约 $0.481$ 个百分点，因此 \textbf{$P_8>0.90$}。结合 $P_n$ 的单调性，可确定
\begin{equation}
 \boxed{N_A^{*}=8}.
\end{equation}
对应的未舍入体积分数为
\[
\Phi_A^{*}=0.0113097336\%.
\]
题目要求体积分数保留百分号后两位，因此对上述精确值舍入得到
\begin{equation}
 \boxed{0.01\%}.
\end{equation}
根数是离散决策变量，因此应先确定 $N_A^*=8$，再对相应体积分数舍入；不能由舍入后的 $0.01\%$ 反推临界根数。

\subsection{结果解释}

本问得到的最低填充量较小，主要与题目规定的边界处理方式有关。A 的长度为 $5\,\mathrm{\mu m}$，等于立方体边长的一半；在各向同性取向下，单根 A 越过 $X$ 向边界的概率约为 $25.47\%$。发生越界后，同一根 A 的两段从立方体相对两侧出现并连接左右电极，因此临界概率主要由单根越界事件决定。$N_A^*=8$ 建立在中心独立均匀、方向各向同性以及题面边界处理方式这些条件上，不应视为一般材料体系的渗流阈值。
% <skill-region:problem-3:verification:start>
\subsection{模型检验：临界根数的独立抽样复核}

问题三的临界整数由概率上下界确定。为复核各向同性方向的抽样方法和单根 A 越过 $X$ 向边界的概率，对轴向采用三维标准高斯向量归一化抽样，并独立生成中心坐标。分别对 $n=7$ 与 $n=8$ 进行 $10^5$ 和 $5\times10^5$ 次抽样，只统计至少一根 A 越过 $X$ 向边界的事件，并将经验频率与解析下界 $L_n$ 比较。解析值应落入相应 Wilson 95\% 区间，且标准化偏差应处于正常随机波动范围，结果见表 \ref{tab:q3-mc-check}。

\begin{table}[H]
\centering
\caption{7 根与 8 根介质 A 越过 $X$ 边界概率的蒙特卡罗复核}
\label{tab:q3-mc-check}
\begin{tabular}{ccccc}
\toprule
$n$ & 样本数 $M$ & 蒙特卡罗估计 & 解析值 $L_n$ & 标准化偏差 $z$ \\
\midrule
7 & $100000$ & $0.872230$ & $0.8722775285$ & $-0.045$ \\
7 & $500000$ & $0.871868$ & $0.8722775285$ & $-0.868$ \\
8 & $100000$ & $0.905730$ & $0.9048100243$ & $0.991$ \\
8 & $500000$ & $0.905210$ & $0.9048100243$ & $0.964$ \\
\bottomrule
\end{tabular}
\end{table}

\textbf{四组解析值均落在对应蒙特卡罗估计的 Wilson 95\% 区间内，且 $|z|<1$}，表明临界点附近的独立随机抽样与解析公式一致。7 根的排除由解析上界 $U_7<0.90$ 保证，8 根的可行性由解析下界 $L_8>0.90$ 保证，因此临界根数的判定不依赖有限样本频率。

\subsection{接触阈值敏感性}

接触阈值 $\delta$ 只影响未越界粒子与电极的接触概率 $p_E=\delta/L$，而单根 A 越过 $X$ 向边界的概率 $p_D$ 与 $\delta$ 无关。为考察临界根数对接触阈值的敏感性，在保持其余条件不变的情况下取
\[
\delta\in\{0,0.9,1.8,3.6,5.4,9.0\}\,\mathrm{nm}.
\]
对每个阈值重新计算 7 根时的导通概率上界。

$U_7$ 从 $0.8722775285$ 缓慢增加至 $0.8722853040$，始终低于 $0.90$。8 根的可行性由 $L_8=0.9048100243$ 保证，该值与 $\delta$ 无关。因此在 $0\le\delta\le9\,\mathrm{nm}$ 的考察范围内，\textbf{$N_A^*=8$ 不变}。

\begin{figure}[H]
\centering
\includegraphics[width=0.82\textwidth]{q7-delta-sensitivity.png}
\caption{接触阈值变化下临界方案的导通概率上界}
\label{fig:q3-delta-sensitivity}
\end{figure}

图 \ref{fig:q3-delta-sensitivity} 同时给出了问题四相关方案的概率上界。对问题三，$U_7$ 在考察区间内始终低于 $0.90$，而与 $\delta$ 无关的 $L_8$ 始终高于 $0.90$，因此题设接触阈值的小范围变化不会改变临界整数。该结论仅适用于 $0\leq\delta\leq9\,\mathrm{nm}$。

\subsection{方向分布敏感性与适用范围}

A 的方向分布会直接改变 $\mathbb E|u_x|$，因而改变单根 A 越过 $X$ 向边界的概率，其影响明显大于接触阈值。经典随机棒体系以及纳米复合材料的蒙特卡罗研究表明，宏观取向或填料对齐会改变渗流阈值与连通性\cite{balberg1984sticks,kale2016alignment}。对平端圆柱，有
\begin{equation}
 p_A^D=\frac{2}{L}\,\mathbb E\!\left[\frac{l_A}{2}|u_x|+r_A\sqrt{1-u_x^2}\right].
 \label{eq:q3-direction-general}
\end{equation}
在完全沿 $X$、各向同性和完全横向三种对照情形下，单根越过 $X$ 向边界的概率分别为 $0.5000000000$、$0.2547123890$ 与 $0.0060000000$；只利用这一事件使概率下界达到 90\% 时，所需 A 根数相应为 4、8 和 383。方向分布改变会明显改变临界填充量，因此 $N_A^*=8$ 只适用于无向轴独立同分布且各向同性的假设。

\begin{figure}[H]
\centering
\includegraphics[width=0.82\textwidth]{q7-direction-sensitivity.png}
\caption{介质 A 方向分布对越过 $X$ 边界概率与临界根数的影响}
\label{fig:q3-direction-sensitivity}
\end{figure}

在所考察范围内，接触阈值变化不会改变临界整数 $8$；介质取向则会明显改变单根越界概率和所需根数。因此，\textbf{问题三对方向分布更敏感}，而对 $0$--$9\,\mathrm{nm}$ 范围内的接触阈值变化不敏感。
% <skill-region:problem-3:verification:end>
% <skill-region:problem-3:end>

% <skill-region:problem-4:start>
\section{问题四：混合介质的最低成本配置}

\subsection{问题分析}

本问在问题三 $90\%$ 导通要求的基础上进一步考虑介质种类和材料成本。\textbf{决策变量为 A、B 两类介质的整数填充数量 $(N_A,N_B)$}；继续采用前述三项随机假设和题面规定的边界处理方式，求导通概率不低于 $90\%$ 时的最低成本配置，并给出两类介质的填充量。多组分填料研究表明，粒子形状、组分比例以及异类粒子之间的连接会影响贯通网络\cite{li2007arbitraryShapes,kale2018multicomponent}。本文仍只采用题设 $\delta=1.8\,\mathrm{nm}$ 的几何接触判据，并按题面规定将越界部分平移到相对侧，不引入隧穿电导模型。

A 较长，单根越过 $X$ 向边界的概率较高；B 的该概率较低，但单位体积价格只有 A 的 $1/21$。只比较单根越界概率、单位体积价格或纯 A、纯 B 两个方案，仍不能排除更便宜的混合配置。完整导通概率 $P_{N_A,N_B}$ 还包含 A--A、A--B 和 B--B 的链式接触。若对所有整数配置逐点进行蒙特卡罗计算，需要反复进行大量几何接触判断；而在 $90\%$ 阈值附近，有限样本只能给出带抽样误差的概率估计，不能作为全局最优性的严格证明。

求解时先用单根介质越过 $X$ 向边界所给出的概率下界证明一个候选配置可行，再以该候选的成本限定所有可能更优的整数配置。若没有粒子越过 $X$ 向边界，任何导电通路都必须分别接触左右电极，据此可以得到更低成本方案的概率上界并逐一排除。由于成本和这一概率上界都随粒子数量单调增加，\textbf{对固定 $N_A$ 只需检查成本限制下允许的最大 $N_B$}。蒙特卡罗抽样只用于复核解析公式和阈值附近的完整网络，不参与最优解的最终判定。

本问先建立成本最小化模型，再用概率下界验证候选方案可行，并对所有成本更低的整数配置逐一排除，最后进行数值复核，整体流程见图 \ref{fig:q4-workflow}。

\begin{figure}[H]
\centering
\includegraphics[width=0.96\textwidth]{q4-workflow.png}
\caption{问题四混合介质最低成本配置的全局优化流程}
\label{fig:q4-workflow}
\end{figure}

\subsection{决策变量、成本函数与基本参数}

设两类介质的填充数量分别为
\[
N_A,N_B\in\mathbb Z_{\ge0}.
\]
介质 A 的尺寸为半径 $r_A=0.03\,\mathrm{\mu m}$、长度 $l_A=5\,\mathrm{\mu m}$，介质 B 的半径为 $r_B=0.2\,\mathrm{\mu m}$。两类单粒子体积分别为
\begin{align}
V_A&=\pi r_A^2l_A=0.01413716694\,\mathrm{\mu m^3},\\
V_B&=\frac43\pi r_B^3=0.03351032164\,\mathrm{\mu m^3}.
\end{align}
题面给出的体积价格分别为 $1.05\,\text{元}/\mu\mathrm{m}^3$ 与 $0.05\,\text{元}/\mu\mathrm{m}^3$，故单根成本为
\begin{align}
c_A&=1.05V_A=0.01484402529\ \text{元},\\
c_B&=0.05V_B=0.001675516082\ \text{元}.
\label{eq:q4-unit-cost}
\end{align}
总成本目标函数为
\begin{equation}
\min C(N_A,N_B)=c_AN_A+c_BN_B,
\label{eq:q4-objective}
\end{equation}
同时满足完整网络导通概率
\begin{equation}
P_{N_A,N_B}\ge0.90.
\label{eq:q4-reliability}
\end{equation}

\subsection{两类介质越过 X 边界的概率与导通概率上下界}

在与问题二、三相同的随机生成假设下，A 单根越过任一 $X$ 向边界的概率为
\begin{equation}
p_A^D=0.254712388980.
\end{equation}
对半径为 $r_B$ 的球体，若球心距离左、右任一 $X$ 边界小于 $r_B$，球体就会越过边界，越出的部分按题面规定从相对侧进入。两个边界对应的球心区间总长度为 $2r_B$，因此
\begin{equation}
p_B^D=\frac{2r_B}{L}=\frac{400}{10000}=0.04.
\label{eq:q4-pb}
\end{equation}
设 $D$ 表示某粒子越过任一 $X$ 边界。由于不同粒子独立，混合配置 $(a,b)$ 中至少有一个粒子发生该事件的概率为
\begin{equation}
L_{a,b}=1-(1-p_A^D)^a(1-p_B^D)^b.
\label{eq:q4-lower}
\end{equation}
只要发生该事件，左右电极即可导通，因此
\[
P_{a,b}\ge L_{a,b}.
\]

为了排除低成本方案，还需给出覆盖所有链式通路的概率上界。对任一未越界粒子，不论其为 A 还是 B，若与指定一侧电极的最短距离不超过接触阈值 $\delta$，其中心在法向方向上的允许区间宽度均为 $\delta$，故与单侧电极接触的概率均为
\begin{equation}
p_E=\frac{\delta}{L}=0.00018.
\end{equation}
若系统没有任何粒子越过 $X$ 向边界，则任何连接左右电极的导电链都必须至少包含一个与左电极接触的粒子和一个与右电极接触的粒子。因此混合配置的导通概率满足
\begin{align}
P_{a,b}\le U_{a,b}
&=1-2(1-p_A^D-p_E)^a(1-p_B^D-p_E)^b \notag\\
&\quad +(1-p_A^D-2p_E)^a(1-p_B^D-2p_E)^b.
\label{eq:q4-upper}
\end{align}
由此得到同一配置的概率上下界
\begin{equation}
L_{a,b}\le P_{a,b}\le U_{a,b}.
\label{eq:q4-bracket}
\end{equation}
其中 $L_{a,b}\ge0.90$ 可以证明该配置可行，$U_{a,b}<0.90$ 可以排除该配置；若 $0.90$ 位于两者之间，仅靠这两个界无法判定，需要进一步计算完整网络。A--A、A--B 和 B--B 链式连接均未被忽略：下界只计入必然导通的一部分情况，上界则覆盖所有导通情况。

\subsection{A、B 单粒子的概率与成本比较}

虽然 A 单根越过 $X$ 边界的概率约为 B 的 $6.37$ 倍，但其单根成本约为 B 的 $8.86$ 倍。以 $p^D/c$ 衡量单根越界概率与成本之比时，A 与 B 分别为 $17.1593$ 与 $23.8732\,\text{元}^{-1}$。若只考虑纯介质方案，由式 \eqref{eq:q4-lower} 得到 A 至少需要 8 根、B 至少需要 57 个，二者成本分别为
\[
C(8,0)=0.1187522023\ \text{元},\qquad
C(0,57)=0.0955044167\ \text{元}.
\]
在给定价格体系下，\textbf{B 每单位成本对应的单根越界概率更高}。这只能说明低成本方案更可能含有较多 B，是否存在更便宜的混合配置仍需检查所有相关整数方案。

\begin{figure}[H]
\centering
\includegraphics[width=0.78\textwidth]{q7-cost-efficiency.png}
\caption{A/B 介质单根越界概率与成本的对比}
\label{fig:q4-cost-efficiency}
\end{figure}

图 \ref{fig:q4-cost-efficiency} 显示，B 的单位成本对应的单粒子越界概率更高，且纯 B 达到 $90\%$ 概率下界时的成本低于纯 A。是否存在成本更低的混合配置，还需继续检查。

\subsection{排除成本更低的整数配置}

由纯 B 方案 $(0,57)$ 的单粒子越界概率下界
\begin{equation}
L_{0,57}=1-0.96^{57}=0.902397648016>0.90,
\label{eq:q4-incumbent-feasible}
\end{equation}
可知\textbf{该方案已经满足 $90\%$ 导通要求}。其成本为
\begin{equation}
\bar C=C(0,57)=0.095504416669\ \text{元}.
\label{eq:q4-cstar}
\end{equation}
若存在成本更低的可行配置，则必须满足 $C(N_A,N_B)<\bar C$。由于
\[
7c_A=0.1039081770\ \text{元}>\bar C,
\]
所以成本低于 $(0,57)$ 的方案只能满足 $0\le N_A\le6$，只需在这 7 个 $N_A$ 取值下考察相应的整数配置。

对固定的 $N_A=a$，定义在严格成本约束下允许的最大 B 数量为
\begin{equation}
 b_a^{\max}
 =\left\lceil\frac{\bar C-ac_A}{c_B}\right\rceil-1,
 \qquad a=0,1,\ldots,6.
 \label{eq:q4-bmax}
\end{equation}
减去 1 用于保证成本严格小于 $\bar C$；当商本身不是整数时，该式与向下取整等价。对固定的 $a$，任何更低成本配置都满足 $0\le b\le b_a^{\max}$。增加 B 粒子不会降低至少一个粒子越过 $X$ 边界、或同时存在与左右电极接触粒子的概率，因此 $U_{a,b}$ 关于 $b$ 单调不减，即
\[
U_{a,b}\le U_{a,b_a^{\max}},\qquad 0\le b\le b_a^{\max}.
\]
因此，对每个固定的 $a$，\textbf{只需检查成本限制下允许的最大 $b$}。若该方案仍满足 $U_{a,b_a^{\max}}<0.90$，则同一 $a$ 下所有更少的 B 也都可以排除。

各 $N_A$ 取值下允许的最大 $N_B$ 及其概率上下界见表 \ref{tab:q4-frontier}。

\begin{table}[H]
\centering
\small
\caption{成本低于 $(0,57)$ 时各 $N_A$ 对应的最大 $N_B$ 及概率界}
\label{tab:q4-frontier}
\begin{tabular}{cccccc}
\toprule
$N_A$ & $N_B^{\max}$ & 成本/元 & $L_{a,b}$ & $U_{a,b}$ & 判定 \\
\midrule
0 & 56 & 0.093828901 & 0.898330883 & 0.898341781 & 不可行 \\
1 & 48 & 0.095268797 & 0.894962812 & 0.894971522 & 不可行 \\
2 & 39 & 0.095033178 & 0.886961628 & 0.886968280 & 不可行 \\
3 & 30 & 0.094797558 & 0.878350957 & 0.878355684 & 不可行 \\
4 & 21 & 0.094561939 & 0.869084369 & 0.869087377 & 不可行 \\
5 & 12 & 0.094326319 & 0.859111902 & 0.859113481 & 不可行 \\
6 & 3  & 0.094090700 & 0.848379784 & 0.848380328 & 不可行 \\
\bottomrule
\end{tabular}
\end{table}

表中七个方案的导通概率上界均低于 $0.90$，其中最接近阈值的是 $(0,56)$，其上界为
\[
U_{0,56}=0.898341781388<0.90.
\]
因此，对任意 $a\in\{0,\ldots,6\}$，即使取成本仍低于 $(0,57)$ 时能够加入的最多 B，导通概率也达不到 $90\%$；同一 $a$ 下更少的 B 更不可能满足要求。另一方面，$N_A\ge7$ 时，仅 A 的成本就已经超过 $\bar C$。于是\textbf{所有成本低于 $(0,57)$ 的非负整数配置都被排除}。

上述七个方案已足以排除全部更低成本配置。再枚举所有满足 $C(N_A,N_B)<\bar C$ 的整数点，共得到 216 个，其导通概率上界均低于 $0.90$，最大值仍为 $(0,56)$ 的 $0.898341781388$。完整枚举与式 \eqref{eq:q4-bmax} 的整数筛选结果一致。

\subsection{全局最优配置与结果解释}

式 \eqref{eq:q4-incumbent-feasible} 已证明 $(0,57)$ 可行；表 \ref{tab:q4-frontier} 排除了所有成本更低的整数配置，216 点完整枚举得到相同结果。因此\textbf{不存在成本小于 $C(0,57)$ 的可行方案}，全局最低成本配置为
\begin{equation}
\boxed{N_A^*=0,\qquad N_B^*=57}.
\label{eq:q4-optimum}
\end{equation}
其最小成本为
\begin{equation}
\boxed{C_{\min}=0.095504416669\ \text{元}}.
\end{equation}
此时介质 A 的填充量为 0；57 个 B 球的总体积为
\[
57V_B=1.9100883334\,\mathrm{\mu m^3},
\]
对应 B 的体积分数为
\begin{equation}
\boxed{\Phi_B=0.191008833338\%}.
\end{equation}

A 更容易越过 $X$ 向边界，但单根成本较高；B 需要更多粒子才能使单粒子越界事件的累计概率超过 $90\%$，较低的价格仍使其总成本低于含 A 的竞争方案。概率下界证明 $(0,57)$ 满足导通要求，概率上界排除了所有更便宜的整数配置。因此，在前述三项随机假设、题面边界处理方式和给定价格条件下，\textbf{$(0,57)$ 是整数范围内的全局最低成本配置。}

% <skill-region:problem-4:verification:start>
\subsection{模型检验：更低成本整数方案的完整枚举复核}

为复核前述排除过程，对所有满足
\[
C(N_A,N_B)<C(0,57)
\]
的非负整数点，共得到 216 个候选。对每个点计算式 \eqref{eq:q4-upper} 的导通概率上界，最大值仍为
\[
U_{0,56}=0.898341781388<0.90.
\]
\textbf{完整枚举所得的最大上界仍低于 $0.90$}，与按每个 $N_A$ 只检查最大 $N_B$ 的解析排除结果一致。

\subsection{几何概率公式的蒙特卡罗复核}

前述证明用到了 A、B 越过 $X$ 向边界的概率，以及未越界粒子与电极接触的概率。对 $(0,57)$ 和前述七个更低成本方案分别进行 $M=10^6$ 次独立抽样，同时估计至少一个粒子越过 $X$ 向边界的频率，以及有粒子越界或同时存在与左右电极接触粒子的频率。A 的 $X$ 向半宽按平端圆柱支撑函数计算，B 按球体几何计算，结果见表 \ref{tab:q4-strict-event-check}。

\begin{table}[H]
\centering
\small
\caption{问题四概率上下界的 $10^6$ 次几何蒙特卡罗复核}
\label{tab:q4-strict-event-check}
\begin{tabular}{ccccc}
\toprule
方案 & MC下界估计 & 解析 $L_{a,b}$ & MC上界估计 & 解析 $U_{a,b}$ \\
\midrule
$(0,57)$ & 0.902561 & 0.902397648 & 0.902577 & 0.902408489 \\
$(0,56)$ & 0.898072 & 0.898330883 & 0.898084 & 0.898341781 \\
$(1,48)$ & 0.894731 & 0.894962812 & 0.894740 & 0.894971522 \\
$(2,39)$ & 0.886512 & 0.886961628 & 0.886515 & 0.886968280 \\
$(3,30)$ & 0.878515 & 0.878350957 & 0.878523 & 0.878355684 \\
$(4,21)$ & 0.869019 & 0.869084369 & 0.869021 & 0.869087377 \\
$(5,12)$ & 0.858746 & 0.859111902 & 0.858746 & 0.859113481 \\
$(6,3)$  & 0.848804 & 0.848379784 & 0.848805 & 0.848380328 \\
\bottomrule
\end{tabular}
\end{table}

\textbf{八个方案的解析下界和上界均落入相应蒙特卡罗估计的 Wilson 95\% 区间。}$(0,57)$ 的解析下界高于 $0.90$，七个更低成本方案的解析上界均低于 $0.90$。独立抽样与解析公式一致，全局最优性仍由前述不等式证明。

\subsection{阈值附近的完整 B 网络独立验证}

为进一步检验阈值附近的完整连通图，选取最优方案 $(0,57)$ 与最接近 $90\%$ 阈值的更低成本方案 $(0,56)$。两者均只含 B 球，可按球--球距离、越界后的平移处理和并查集连通判定进行完整网络蒙特卡罗验证。每个方案采用 3 个独立随机种子，每次抽样 $M=10^6$，结果见表 \ref{tab:q4-b-only-check}。

\begin{table}[H]
\centering
\small
\caption{纯 B 临界方案的三随机种子完整连通图验证}
\label{tab:q4-b-only-check}
\begin{tabular}{ccccc}
\toprule
方案 & 随机种子 & $\hat p$ & Wilson 95\% 区间 & 仅链式网络导通频率 \\
\midrule
$(0,57)$ & 2026081801 & 0.902176 & $[0.901592,\,0.902757]$ & 0 \\
$(0,57)$ & 2026081807 & 0.902288 & $[0.901704,\,0.902868]$ & 0 \\
$(0,57)$ & 2026081819 & 0.902719 & $[0.902137,\,0.903298]$ & 0 \\
$(0,56)$ & 2026081801 & 0.898145 & $[0.897551,\,0.898736]$ & 0 \\
$(0,56)$ & 2026081807 & 0.898184 & $[0.897590,\,0.898775]$ & 0 \\
$(0,56)$ & 2026081819 & 0.898611 & $[0.898018,\,0.899201]$ & 0 \\
\bottomrule
\end{tabular}
\end{table}

\textbf{$(0,57)$ 三次经验概率的均值为 $0.902394333$，且三次 Wilson 区间下端点均高于 $0.90$；$(0,56)$ 三次均值为 $0.898313333$，且三次 Wilson 区间上端点均低于 $0.90$。}六次试验中仅由链式网络产生的导通频率均为 0，说明在当前纯 B 临界填充范围内，导通主要由单个 B 球越过 $X$ 边界产生。该结论不能外推到更高填充量，因为 B--B 链式通路的出现概率可能随粒子数量增加而上升。

\subsection{接触阈值扰动与适用范围}

为检验最优整数配置对接触阈值的稳定性，保持粒子分布和价格不变，将 $\delta$ 从 $0$ 扰动至 $9\,\mathrm{nm}$。候选 $(0,57)$ 的可行性下界
\[
L_{0,57}=0.902397648016
\]
仅由单个粒子越过 $X$ 边界的概率决定，因此在该区间内保持不变；所有更便宜方案中，导通概率上界最高的始终是 $(0,56)$，其上界由 $0.898330883350$ 增至 $0.898592567280$，仍低于 $0.90$。在 $0\le\delta\le9\,\mathrm{nm}$ 的考察范围内，该上界始终低于 $0.90$，因此 \textbf{$(0,57)$ 仍为全局最低成本配置}，见图 \ref{fig:q3-delta-sensitivity}。

介质 A 的方向分布会明显改变 $p_A^D$。若 A 从各向同性变为强烈沿 $X$ 取向，部分含 A 的低成本混合方案可能具有更高的单根越界概率，因此需要重新比较所有相关整数配置。$(0,57)$ 在所考察的接触阈值范围内保持最优，但这一结论不能直接用于其他方向分布。

对于包含 A 的混合体系，精确计算完整网络概率时应采用平端圆柱--圆柱及圆柱--球体的实体距离。本文的全局最优性证明不依赖 A--A、A--B 链式距离，且最优方案和最接近阈值的更低成本方案均为纯 B，因此这一几何细化不影响本题的最低成本结论；若进一步研究混合网络的精确导通概率或成本与导通概率的多目标关系，则需重新计算完整实体距离。
% <skill-region:problem-4:verification:end>
% <skill-region:problem-4:end>

% <skill-region:synthesis-evaluation:start>

\section{模型评价与推广}

\subsection{模型优点}

（1）\textbf{四个问题采用一致的边界处理方式。}介质越出立方体后，将越出的部分平移一个边长，从相对一侧继续参与接触判断。问题一据此判断附件中位于相对边界的两段记录是否属于同一根介质，问题二至问题四则在相同规则下计算单根介质连接两侧端面的概率。四问分别处理确定性数据、随机填充和成本优化，但导通判定使用同一套几何规则。

（2）\textbf{关键填充阈值可以由解析概率上下界直接判断。}问题二将至少有一根介质自身连接左右两个端面作为导通的充分条件，由此得到四种体积分数下的概率下界。问题三在此基础上，再利用导通所必须满足的端面接触条件构造概率上界，并由 $U_7<0.90$ 与 $L_8>0.90$ 分别排除 7 根、确认 8 根达到要求。蒙特卡罗抽样只用于复核这些几何概率，没有替代解析判定。

（3）\textbf{问题四用概率上下界确定最低成本方案。}本文先证明 57 个 B 球组成的方案满足 90\% 的导通要求；随后对 $N_A=0,1,\ldots,6$，分别取成本低于 $(0,57)$ 时允许的最大 $N_B$，计算相应的导通概率上界。七个上界均低于 0.90，而同一 $N_A$ 下减少 $N_B$ 只会使该上界进一步降低，因此所有更低成本方案均可排除。对 216 个更低成本整数配置的完整枚举得到相同结论。

（4）\textbf{灵敏度计算给出了结论成立的范围。}在 $0\leq\delta\leq9\,\mathrm{nm}$ 内，问题三的临界根数仍为 8，问题四的最低成本配置仍为 $(0,57)$；A 的方向分布则会明显改变其连接两侧端面的概率，并可能改变临界根数和混合方案的成本比较。因此，接触阈值在当前考察范围内对最终答案影响较小，而介质取向是需要重点确认的假设。

\subsection{模型局限}

（1）\textbf{问题二至问题四采用的随机分布没有实测数据标定。}题目未给出介质中心位置和方向的统计规律，本文假设介质中心相互独立且均匀分布，A 的轴向各向同性。若实际制备过程中存在明显取向、团聚、排斥或空间浓度梯度，单根介质连接两侧端面的概率会发生变化，临界填充量和最低成本配置也需要重新计算。

（2）\textbf{解析上下界用于判断本题给出的阈值，但不能给出任意填充量下的精确导通概率。}上下界之间还包含由 A--A、A--B 和 B--B 多粒子接触形成的导电路径。若研究目标改为得到完整的填充量--导通概率曲线，就需要显式生成所有介质并计算它们之间的接触关系。

（3）\textbf{问题一只能利用附件中能够严格核对的记录。}本文要求相对 $X$ 边界处的两段记录同时满足端点位置对应、方向共线和长度之和等于 $5000\,\mathrm{nm}$，据此已能确认三组样本均存在连接左右端面的 A。附件中其余不满足这些条件的短段和异常记录没有被补配，因此当前结果足以回答是否导通，但不能据此恢复三组样本的全部介质结构。

（4）\textbf{若要精确比较含 A 的混合网络，A--A 和 A--B 接触距离还应按平端圆柱实体计算。}本文问题四的最低成本证明只使用单粒子越过边界的解析概率和纯 B 临界方案，因此不受这一点影响；但在进一步研究含 A 方案的完整导通概率时，需要替换现有的圆柱轴线近似并重新验证数值结果。

\subsection{改进方向与适用范围}

\textbf{若能获得显微成像、取向统计或制备工艺数据}，可直接估计介质中心位置的空间分布、位置相关性和 A 的方向分布，并将这些统计量代入随机几何模型。这样可以把问题二至问题四中的随机假设改为由实际材料数据确定的分布。

\textbf{对于需要精确导通概率的混合体系}，可统一采用平端圆柱--圆柱、圆柱--球体和球体--球体的实体距离，逐次生成介质接触网络并判断两个端面是否连通。在概率接近目标阈值时，还可采用分层抽样、准蒙特卡罗或重要抽样减少估计方差，并同时给出置信区间。

\textbf{本文方法适用于立方体区域中、介质越出一个端面后从相对端面继续出现、且导通由几何接触决定的填充问题。}若改变粒子形状、A 的取向分布、材料单价或边界条件，需要重新计算单粒子的几何概率及单位成本；若边界改为直接截断，或允许跨越相对边界的独立介质直接接触，则应先重新定义接触和导通判据，再求相应的可靠性阈值和最低成本方案。
% <skill-region:synthesis-evaluation:end>

\bibliographystyle{gbt7714-numerical}
\bibliography{references}

\appendix

% <skill-region:appendix-code:start>
\section{核心代码}
\label{app:core-code}

\subsection{问题四：最低成本整数解的解析判定}
\lstinputlisting[
  firstline=4,lastline=38,
  caption={低于候选最优成本的整数方案概率上界计算},
  label={lst:q4-global-proof}
]{03_code/q4/q4_global_proof.py}

\subsection{问题四：越界处理与介质接触网络计算}
\lstinputlisting[
  linerange={82-92,94-103,105-123,125-143},
  caption={A、B 介质越出立方体后的对侧平移与位置记录},
  label={lst:q4-boundary-kernel}
]{03_code/q4/q4_strict_periodic_mc.py}

\lstinputlisting[
  firstline=217,lastline=345,
  caption={并查集构造介质接触网络并判断左右端面是否连通},
  label={lst:q4-mc-kernel}
]{03_code/q4/q4_strict_periodic_mc.py}

\subsection{问题三与问题四：接触阈值敏感性计算}
\lstinputlisting[
  linerange={45-49,51-107,109-113},
  caption={概率上下界与接触阈值敏感性计算},
  label={lst:q7-sensitivity}
]{03_code/analysis/q7_sensitivity_analysis.py}
% <skill-region:appendix-code:end>

% <skill-region:appendix-support-list:start>
\section{支撑材料列表}
\label{app:support-files}

表~\ref{tab:support-files}列出与正文主要结果直接对应的原始数据、程序和计算结果文件。

{\footnotesize
\begin{longtable}{>{\raggedright\arraybackslash}p{0.23\textwidth}>{\raggedright\arraybackslash}p{0.08\textwidth}>{\raggedright\arraybackslash}p{0.28\textwidth}>{\raggedright\arraybackslash}p{0.21\textwidth}}
\caption{主要支撑材料及对应结果}\label{tab:support-files}\\
\toprule
文件名 & 所属问题 & 主要功能 & 对应结果 \\
\midrule
\endfirsthead
\multicolumn{4}{c}{续表~\thetable}\\
\toprule
文件名 & 所属问题 & 主要功能 & 对应结果 \\
\midrule
\endhead
\midrule
\multicolumn{4}{r}{下页续}\\
\endfoot
\bottomrule
\endlastfoot
附件.xlsx & 问题一 & 提供三组介质 A 的轴线端点原始数据 & 三组样本的导通判定 \\
\path{q1_complete_x_components_verified.csv} & 问题一 & 核对相对 $X$ 边界记录的端点位置、方向和长度 & 符合三项条件的记录对为 2、5、63 \\
\path{q4_global_proof.py} & 问题四 & 计算成本低于 $(0,57)$ 的整数方案的导通概率上界 & $(0,57)$ 为最低成本配置 \\
\path{stage3_all_216_cheaper_points.csv} & 问题四 & 枚举成本低于 $(0,57)$ 的全部 216 个整数配置 & 216 个配置均不能达到 90\% 要求 \\
\path{stage4_strict_event_M1000000.csv} & 问题四 & 复核 A、B 单粒子越过 $X$ 边界及接触端面的几何概率 & 8 个关键配置的模拟结果与解析计算一致 \\
\path{stage5_final_B_only_M1000000.csv} & 问题四 & 对 56B 和 57B 的完整接触网络进行独立抽样 & 57B 与 56B 分别位于 0.90 两侧 \\
\path{q7_sensitivity_analysis.py} & 问题三、四 & 计算接触阈值和 A 的方向分布改变时的最小 A 根数与最低成本配置 & $N_A=8$ 与 $(0,57)$ 在 $0$--$9$ nm 内不变 \\
\path{q4_strict_periodic_mc.py} & 问题四 & 生成越界后的对应位置，计算介质接触并判断左右端面连通 & 56B 与 57B 配置的独立验证 \\
\end{longtable}
}
% <skill-region:appendix-support-list:end>

\end{document}
